\documentclass[10pt,a4paper]{amsart}
\usepackage[margin=19mm]{geometry}
\usepackage{amsmath,amssymb,mathtools}
\usepackage[scr=rsfs,cal=euler]{mathalfa}
\usepackage{xfrac}
\usepackage[unicode,hidelinks,bookmarks=false]{hyperref}
\newcommand{\ie}{i.e.\ }
\newcommand{\cf}{cf.\ }
\newcommand{\resp}{respectively\ }
\newcommand{\Subsection}{\subsection}
\providecommand{\nobreakdash}{\nobreak\hbox{-}}
\setlength{\emergencystretch}{2em}
\multlinegap0pt
\usepackage{amssymb}
\usepackage{mathtools}
\usepackage{tikz}
\usepackage{csquotes}
\newtheorem{prop}{Proposition}
\newtheorem{lemma}{Lemma}
\def\Hom{H^1_0(\Omega)}
\def\Sobom#1{W^{1,#1}_0(\Omega)}
\newcommand{\abs}[1]{\left|#1\right|}
\def\elle#1{L^{#1}}
\def\elleom#1{L^{#1}(\Omega)}
\def\elleomT#1{L^{#1}(\Omega_T)}
\def\io{\int_\Omega}
\def\ioT{\int_{\Omega_T}}
\def\iot{\int_{\Omega_t}}
\def\norma#1#2{ \|#1 \|_{#2}}
\title{Existence and regularity of solutions to parabolic-elliptic nonlinear systems}
\author{Marco Picerni}
\date{Source: arXiv:2604.16100v2 (CC BY 4.0)}
\begin{document}
\maketitle

\noindent\emph{Source and licence.} Existence and regularity of solutions to parabolic-elliptic nonlinear systems. Version arXiv:2604.16100v2 (CC BY 4.0), distributed by arXiv under the Creative Commons Attribution 4.0 International licence, http://creativecommons.org/licenses/by/4.0/, as recorded in the arXiv licence metadata for this exact version and shown on the abstract page https://arxiv.org/abs/2604.16100v2. Full licence notice: \texttt{licenses/arxiv-2604.16100-CC-BY-4.0.txt}.\par\smallskip

\noindent\textit{Editorial scope.} Counted unit: the article's lemma lemma:stima1** (source0.tex lines 981--983). The article's complete original proof of it is delivered with it (source0.tex lines 984--1153, one \texttt{proof} environment of 170 lines). No mathematics is written, repaired or re-derived: the statement and the proof are the article's own lines, in the article's own notation, and no retained line is substituted or re-flowed.  Retained uncounted as the source-local context the proof needs: lines 241--249; lines 400--408; lines 478--481; lines 810--812; lines 873--875.  Nothing was omitted from any retained span, so the package carries no omission record.  No retained line contains a citation, so the wrapper carries no bibliography and no bibliography entry.  No retained line contains a figure environment, an image-inclusion command or drawing code, so no image is delivered and the package declares no asset dependency.  Editorial formatting only, in addition: an AMS \texttt{amsart} preamble that restates the article's own declarations and macros used by the retained lines, namely the environments \texttt{prop}, \texttt{lemma} on separate counters, together with the standard packages those lines need.  The retained environments print in this document as Proposition 1, Lemma 1 in order of appearance, whereas the article prints the same items with the numbering of its own full document.  The complete native source of the article as distributed in the arXiv e-print for this exact version is bundled unchanged as \texttt{sources/198664/source0.tex}.
\begin{equation}\label{eq:approx_KS_system}
    \left\{\begin{array}{cc}
    (u_n)_t-\operatorname{div}(A(x, t) \nabla u_n)=-\operatorname{div}(T_n(u_n) M(x) \nabla \psi_n)+T_n(f(x,t)) & \text { in } \Omega_T, \\
    -\operatorname{div}(M(x) \nabla \psi_n)=|T_n(u_n)|^\theta & \text { in } \Omega_T, \\
    \psi_n(x, t)=0 & \text { on } \partial\Omega\times(0, T), \\
    u_n(x, t)=0 & \text { on } \partial\Omega\times(0, T), \\
    u_n(x, 0)=0 & \text { in } \Omega.
    \end{array}\right.
\end{equation}

\begin{prop}\label{GNSparabTheo}
    Let $v\in\elle\infty(0,T,\elleom2)\cap\elle2(0,T;\Hom)$. Then $v\in\elle{2\frac{N+2}{N}}(\Omega_T)$  and
    \begin{equation}\label{GNSparabEQ}
        \ioT \abs{v}^{2\frac{N+2}{N}}
        \leq C
        \norma{v}{\elle\infty(0,T;\elleom2)}^\frac{4}{N}\norma{\nabla v}{\elle2(\Omega_T)}^2,      
    \end{equation}
    where $C$ is a positive constant which depends on $N$ and $\Omega$.
\end{prop}

\begin{lemma}\label{StimaL1ParabEll}
Let $u_n$ be a solution to the first equation to \eqref{eq:approx_KS_system}, then $u_n\in\elle\infty(0,T;\elleom1)$ and
$$\norma{u_n}{\elle\infty(0,T;\elleom1)}\leq \norma{f}{\elle1(\Omega_T)}.$$
\end{lemma}

\begin{lemma}\label{lemma:stima_m**_unif:infinite_energy}
    Let $f\in\elle m (\Omega_T)$ with $1<m<\frac{2N+4}{N+4}$, then the sequence $(u_n)_n$ is bounded in $\elleomT {m^{\star\star}}$.
\end{lemma}

\begin{lemma}\label{lemma:u_nW1m*_KS_parabell}
    Let $f\in\elleomT m$, with $1<m<\frac{2N+4}{N+4}$. Then the sequence of approximating solutions is bounded in $\elle {m^\star}(0,T;\Sobom{m^\star})$ (where ${m^\star}=\frac{(N+2)m}{(N+2)-m})$.
\end{lemma}

\begin{lemma}\label{lemma:stima1**}
    Let $f\in\elleomT1$, then the sequence of approximating solutions $(u_n)_n$ is bounded in $\elleomT p$ for every $p<\frac{N+2}{N}=1^{\star\star}$ and the sequence of gradients $(\nabla u_n)_n$ is bounded in $\elleomT q$ for every $q<\frac{N+2}{N+1}=1^{\star}$.
\end{lemma}

\begin{proof}
    Let $1<\lambda<2$ and choose $\frac{1-(1+u_n)^{1-\lambda}}{\lambda-1}$ as a test function in the second equation of \eqref{eq:approx_KS_system}. Repeating the same steps as in the proof of Lemma \ref{lemma:stima_m**_unif:infinite_energy}, we define the function
    $$F_n(s)=\int_0^s (\sigma+1)^{-\lambda}T_n(\sigma) \mathrm{d}\sigma$$
    and choose $F_n(u_n)$ as a test function in the second equation of \eqref{eq:approx_KS_system} to get
    $$
    \io (u_n+1)^{-\lambda}T_n(u_n) M(x)\nabla \psi_n\cdot\nabla u_n 
    =
    \io M(x)\nabla \psi_n\cdot\nabla F_n(u_n)
    =
    \io T_n(u_n)^\theta F_n(u_n)
    $$
    $$
    \leq \io T_n(u_n)^\theta \frac{(u_n+1)^{2-\lambda}-1}{2-\lambda}
    \leq
    \io u_n^{\theta} \frac{(u_n+1)^{2-\lambda}-1}{2-\lambda}
    $$
    Now choose 
    $$g(u_n)=\frac{1-(1+u_n)^{1-\lambda}}{\lambda-1}$$
    (which is bounded from above by $\frac{1}{\lambda-1}$) as a test function in the first equation of \eqref{eq:approx_KS_system} to get
\begin{equation}\label{eq:proof:stima1**:eq0}
    \begin{split}
            \langle (u_n)_t,g(u_n)\rangle&
    +\io(u_n+1)^{-\lambda}A(x,t)\nabla u_n \cdot \nabla u_n\\&
    \leq
    \io (u_n+1)^{-\lambda}T_n(u_n) M(x)\nabla \psi_n\cdot\nabla u_n 
    +\io f(x,t)g(u_n).
    \end{split}
\end{equation}
    Which implies
    \begin{equation}\label{EQ_stimagrad_Parabell_fL1}
        \langle (u_n)_t,g(u_n)\rangle+\alpha\io\frac{\abs{\nabla u_n}^2}{(u_n+1)^{\lambda}}
        \leq
        \io u_n^{\theta} \frac{(u_n+1)^{2-\lambda}-1}{2-\lambda}+\frac{1}{\lambda-1}\io f(x,t).        
    \end{equation}
    We now focus on the term
    $$\io u_n^{\theta} \frac{(u_n+1)^{2-\lambda}-1}{2-\lambda}.$$
    Let $k\geq0$ and apply the H\"older inequality (with exponent $\frac{N}{2}$) along with Lemma \ref{StimaL1ParabEll} to get
    \begin{equation}\label{eq:proof:stima1**:eq1}
        \begin{split}
            \io u_n^{\theta} &\frac{(u_n+1)^{2-\lambda}-1}{2-\lambda}
    \leq \int_{u_n(t)\leq k} u_n^{\theta} \frac{(u_n+1)^{2-\lambda}-1}{2-\lambda} + 
    \int_{u_n(t)> k} u_n^{\theta} \frac{(u_n+1)^{2-\lambda}-1}{2-\lambda}\\
    &\leq
    k^{\theta}\frac{(k+1)^{2-\lambda}}{2-\lambda}\abs{\Omega}+\frac{1}{2-\lambda}
    \left(\frac{\norma{f}{\elleomT1}}{k^{1-\theta\frac{N}{2}}}\right)^\frac{2}{N}
    \left(\io \left((u_n(t)+1)^{2-\lambda}-1\right)^{\frac{N}{N-2}}\right)^{\frac{N-2}{N}}.
        \end{split}
    \end{equation}
    
    Since, for every $\eta\geq0$, the following inequality holds
    $$
    \left[(1+\eta)^{2-\lambda}-1\right]^{\frac{1}{2}} \leq(1+\eta)^{\frac{2-\lambda}{2}}=\left[(1+\eta)^{\frac{2-\lambda}{2}}-1\right]+1,
    $$
    we have
    $$
    \io\left[\left(1+u_n\right)^{2-\lambda}-1\right]^{\frac{N}{N-2}} 
    \leq 
    \io\left\{\left[\left(1+u_n\right)^{\frac{2-\lambda}{2}}-1\right]+1\right\}^{{\frac{2N}{N-2}}} 
    \leq 
    C \left(1+\io\left[\left(1+u_n\right)^{\frac{2-\lambda}{2}}-1\right]^{{\frac{2N}{N-2}}}\right),
    $$
    for some $C>0$.
    It follows that
        \begin{equation}\label{eq:proof:stima1**:eq3}
        \begin{split}
            \io u_n^{\theta} \left((u_n+1)^{2-\lambda}-1\right)
    \leq 
    k^{\theta+2-\lambda}\abs{\Omega}
    +    \frac{C}{2-\lambda}
    \left(\frac{\norma{f}{\elleomT1}}{k^{1-\theta\frac{N}{2}}}\right)^\frac{2}{N}
    \left(1+\left[\io\left[\left(1+u_n\right)^{\frac{2-\lambda}{2}}-1\right]^{{\frac{2N}{N-2}}}\right]^{\frac{N-2}{N}}\right).
        \end{split}
    \end{equation}

    On the other hand, by the Sobolev inequality, we have
    \begin{equation}\label{eq:proof:stima1**:eq2}
        \begin{split}
            \io \frac{\left|\nabla u_n\right|^2}{\left(1+u_n\right)^\lambda}=
    \frac{4}{(2-\lambda)^2} \int_{\Omega}\left|\nabla\left[\left(1+u_n\right)^{\frac{2-\lambda}{2}}-1\right]\right|^2 \geq 
    \frac{4 \mathcal{S}}{(2-\lambda)^2}\left[\int_{\Omega}\left[\left(1+u_n\right)^{\frac{2-\lambda}{2}}-1\right]^{{\frac{2N}{N-2}}}\right]^{\frac{N-2}{N}}.
        \end{split}
    \end{equation}
    
    Notice that, if we define $\mathcal{G}(s)=\int_0^s g(\sigma)\mathrm{d}\sigma$, the term on the left-hand side of \eqref{EQ_stimagrad_Parabell_fL1} becomes $\frac{d}{dt}\io \mathcal{G}(u_n)$. Thus, putting together \eqref{EQ_stimagrad_Parabell_fL1}, \eqref{eq:proof:stima1**:eq3} and \eqref{eq:proof:stima1**:eq2}, we obtain
    \begin{equation}
    \begin{split}
    \frac{d}{dt}\io \mathcal{G}(u_n) +
    &\left(\alpha -C \frac{(2-\lambda)^2}{4\mathcal{S}}\frac{\norma{f}{\elleomT1}^{\frac{2}{N}}}{(2-\lambda) k^{\frac{2}{N}-\theta}}\right)
    \io \frac{\left|\nabla u_n\right|^2}{\left(1+u_n\right)^\lambda} \\
    & \leq 
    \frac{1}{\lambda-1}\io f(x,t) +
    k^{\theta}\frac{(k+1)^{2-\lambda}}{2-\lambda}\abs{\Omega}
    +\frac{C}{2-\lambda} \frac{\norma{f}{\elleomT 1} ^{\frac{2}{N}}}{k^{\frac{2}{N}-\theta}} .
    \end{split}
    \end{equation}
    Choosing $k=k_0$ such that 
    $$\alpha -C \frac{(2-\lambda)^2}{4\mathcal{S}}\frac{\norma{f}{\elleomT1}^{\frac{2}{N}}}{(2-\lambda) k_0^{\frac{2}{N}-\theta}}=\frac{\alpha}{2},$$
    we have
    $$\frac{d}{dt}\io \mathcal{G}(u_n) +
    \frac{\alpha}{2}\io \frac{\left|\nabla u_n\right|^2}{\left(1+u_n\right)^\lambda}
    \leq 
    \frac{1}{\lambda-1}\io f(x,t) +
    k_0^{\theta}\frac{(k_0+1)^{2-\lambda}}{2-\lambda}\abs{\Omega}
    +\frac{C}{2-\lambda} \frac{\norma{f}{\elleomT 1} ^{\frac{2}{N}}}{k_0^{\frac{2}{N}-\theta}} .
    $$
    Integrating from $0$ to $t$, we obtain
    $$\io \mathcal{G}(u_n)(t) +
    \frac{\alpha}{2}\iot \frac{\left|\nabla u_n\right|^2}{\left(1+u_n\right)^\lambda}
    \leq 
    \frac{1}{\lambda-1}\ioT f(x,t) +
    k_0^{\theta}\frac{(k_0+1)^{2-\lambda}}{2-\lambda}\abs{\Omega}T
    +\frac{C}{2-\lambda} \frac{\norma{f}{\elleomT 1} ^{\frac{2}{N}}}{k_0^{\frac{2}{N}-\theta}}T,
    $$
    which is
    \begin{equation}\label{eq:proof:stima1**:eq3.5}
        \io \mathcal{G}(u_n)(t) +
    \frac{\alpha}{2}\iot \frac{\left|\nabla u_n\right|^2}{\left(1+u_n\right)^\lambda}
    \leq C\left(\|f\|_{L^1(\Omega_T)}, \lambda\right),
    \end{equation}
    where $C\left(\|f\|_{L^1(\Omega_T)},\lambda\right)$ is a positive constant which tends to infinity as $\lambda\to1^+$.

    Since $\mathcal{G}(s)=s+(1+s)^{2-\lambda}-1$ (which, for $s\geq0$, is greater than $s^{2-\lambda}$) and $u_n\geq0$, we have, up to increasing the constant $C\left(\|f\|_{L^1(\Omega_T)},\lambda\right)$,
    \begin{equation}\label{DoppiaStima-Lemma-StimaconfL1}
        \left\{\begin{aligned}
        & \norma{(u_n+1)^{2-\lambda}-1}{\elle\infty(0,T;\elleom1)}\leq C\left(\|f\|_{L^1(\Omega_T)}, \lambda\right)\\
        & \norma{\nabla \left((u_n+1)^\frac{2-\lambda}{2}-1\right)}{\elleomT2}^2\leq C\left(\|f\|_{L^1(\Omega_T)}, \lambda\right).
        \end{aligned}\right.
    \end{equation}
    Recall that, by Proposition \eqref{GNSparabTheo}, there exists $C=C(N,\Omega)>0$ such that
    $$\norma{(u_n+1)^\frac{2-\lambda}{2}-1}{\elleomT{2\frac{N+2}{N}}}
    \leq\norma{(u_n+1)^\frac{2-\lambda}{2}-1}{\elle\infty(0,T;\elleom2)}^\frac{4}{N}
    \norma{\nabla \left((u_n+1)^\frac{2-\lambda}{2}-1\right)}{\elleomT2}^2.
    $$
    Moreover, since $\frac{2-\lambda}{2}\in (0,1)$ and $u_n\geq0$ we have, by concavity,
    $$(u_n+1)^\frac{2-\lambda}{2}-1\leq u_n^\frac{2-\lambda}{2},$$
    which implies
    $$
    \norma{(u_n+1)^\frac{2-\lambda}{2}-1}{\elle\infty(0,T;\elleom2)}
    \leq 
    \norma{u_n^\frac{2-\lambda}{2}}{\elle\infty(0,T;\elleom2)}
    \leq
    \norma{u_n^{2-\lambda}}{\elle\infty(0,T;\elleom1)}^\frac{1}{2}.$$
    By convexity, since $2-\lambda>1$, we get
    $$\norma{(u_n+1)^\frac{2-\lambda}{2}-1}{\elle\infty(0,T;\elleom2)}
    \leq 
    \norma{u_n^{2-\lambda}}{\elle\infty(0,T;\elleom1)}^\frac{1}{2}
    \leq
    \norma{(u_n+1)^{2-\lambda}-1}{\elle\infty(0,T;\elleom1)}^\frac{1}{2}.$$
    To sum up, we have proved that, for any $\lambda\in(1,2)$, there exists a positive constant $C=C\left(\norma{f}{\elleomT1},\lambda,N,\Omega\right)$ such that
    \begin{equation}\label{eq:proof:stima1**:eq4}
        \norma{(u_n+1)^\frac{2-\lambda}{2}-1}{\elleomT{2\frac{N+2}{N}}}
    \leq C,
    \end{equation}
    which implies that the sequence $(u_n^\frac{2-\lambda}{2})_n$ is bounded in $\elleomT{2\frac{N+2}{N}}$. This means that $(u_n)_n$ is bounded in $\elleomT{{2-\lambda}\frac{N+2}{N}}$ for any $\lambda\in(1,2)$, i.e. $(u_n)_n$ is bounded in $\elleomT p$ for any $p<\frac{N+2}{N}$.

    We now follow the argument in the proof of Lemma \ref{lemma:u_nW1m*_KS_parabell}. Choosing $q<2$ and using the H\"older inequality, we obtain 
    $$\ioT \abs{\nabla u_n}^q =
    \ioT \frac{\abs{\nabla u_n}^q}{(1+u_n)^{\frac{q\lambda}{2}}}(1+u_n)^{\frac{q\lambda}{2}} 
    \leq
    \left(\ioT \frac{\abs{\nabla u_n}^2}{(1+u_n)^{\lambda}}\right)^\frac{q}{2}
    \left(\ioT (1+u_n)^{\frac{q\lambda}{2-q}}\right)^{1-\frac{q}{2}},$$
    thus, by \eqref{eq:proof:stima1**:eq3.5}, we get
    $$\ioT \abs{\nabla u_n}^q\leq 
     \frac{2\,C\left(\norma{f}{\elleomT1}, \lambda\right)}{\alpha}
    \left(\ioT (1+u_n)^{\frac{q\lambda}{2-q}}\right)^{1-\frac{q}{2}}.$$

    Setting $\frac{q\lambda}{2-q}={2-\lambda}\frac{N+2}{N}$ and recalling \eqref{eq:proof:stima1**:eq4}, we conclude that $(\nabla u_n)_n$ is bounded in $\elleomT q$ with
    $$q=q(\lambda)=\frac{2(2-\lambda)(N+2)}{\lambda N +(2-\lambda)(N+2)}.$$
    Note that $q(\lambda)$ is always smaller that $\frac{N+2}{N+1}$ and tends to $\frac{N+2}{N+1}$ as $\lambda\to1$. This concludes the proof.
\end{proof}

\end{document}
